% ===========================================================================
\section{Proposed System Design}
% Source material: docs/architecture.md sections 3-7, docs/tech-stack.md
% ===========================================================================

\subsection{Architecture overview}

\projectname{} has three layers (Figure~\ref{fig:architecture}): a Next.js web
app with the admin console, the ballot booth and the audit dashboard; four Node
services that have no say over the result (registrar, relayer, indexer, IPFS
metadata); and the smart contracts, which hold everything that decides the
election.

\begin{figure}[htbp]
  \centering
  % DVote system architecture.
% Bare tikzpicture: \input this inside a figure environment.
% The preamble supplies tikz with arrows.meta, positioning, shapes.geometric,
% fit, backgrounds and calc.
%
% Reading conventions, all greyscale-safe:
%   rounded outline = browser surface       dashed outline = off-chain service
%   heavy outline   = on-chain contract     cylinder       = persistent store
%   thin grey arrow = ordinary traffic
%   heavy solid     = the sponsored ballot path (voter -> relayer -> chain)
%   heavy dashed    = the audit path, which never touches the backend
%   fine dotted     = the escape hatch, a voter submitting their own transaction
%
% Every off-chain writer that reaches the chain has its own arrow: the relayer
% into VoteForwarder, the registrar into VoterRegistry, the authority's own
% wallet into Election. Nothing in the middle band decides the result.
\begin{tikzpicture}[
    font=\footnotesize,
    uibox/.style    = {draw=black, thick, rounded corners=6pt, fill=white,
                       align=center, inner sep=3pt},
    svcbox/.style   = {draw=black, thick, dashed, rounded corners=2pt, fill=white,
                       align=center, inner sep=3pt},
    store/.style    = {cylinder, shape border rotate=90, aspect=0.20, draw=black,
                       thick, fill=white, align=center, inner sep=3pt},
    cbox/.style     = {draw=black, line width=1.1pt, fill=black!7,
                       align=center, inner sep=3pt},
    layer/.style    = {draw=black!30, rounded corners=5pt, fill=black!4, inner sep=0pt},
    layerlab/.style = {font=\footnotesize\itshape, text=black!60, inner sep=2pt},
    flow/.style     = {-{Latex[length=2mm,width=1.6mm]}, draw=black!60, semithick},
    ballot/.style   = {-{Latex[length=2.6mm,width=2.2mm]}, draw=blue!45!black,
                       line width=1.3pt, rounded corners=4pt},
    audit/.style    = {-{Latex[length=2.6mm,width=2.2mm]}, draw=blue!45!black,
                       line width=1.3pt, rounded corners=4pt,
                       dash pattern=on 4.5pt off 2.5pt},
    hatch/.style    = {-{Latex[length=2mm,width=1.6mm]}, draw=black!70, semithick,
                       rounded corners=4pt, dash pattern=on 1pt off 1.7pt},
    lab/.style      = {align=left, inner sep=1.5pt, font=\footnotesize},
    keylab/.style   = {align=left, inner sep=1.5pt, font=\footnotesize,
                       text=blue!45!black},
    slab/.style     = {align=center, inner sep=1.5pt, font=\scriptsize,
                       fill=black!4, text=black!70},
  ]

  % ==== browser layer ====
  \node[uibox, minimum width=4.6cm, minimum height=1.5cm] (booth) at (3.75,-0.55)
    {\textbf{Ballot booth}\\rebuilds the tree from the\\published leaf set, proves\\in a Web Worker, no wallet};
  \node[uibox, minimum width=4.4cm, minimum height=1.5cm] (admin) at (9.2,-0.55)
    {\textbf{Admin console}\\configure, open each phase,\\cancel before voting;\\the authority wallet signs};
  \node[uibox, minimum width=4.4cm, minimum height=1.5cm] (audit) at (14.1,-0.55)
    {\textbf{Audit dashboard}\\tally recomputed from raw\\logs in the browser, over\\a public RPC endpoint};

  % ==== backend layer ====
  \node[svcbox, minimum width=2.8cm, minimum height=1.9cm] (relayer) at (3.4,-4.8)
    {\textbf{Relayer}\\verifies the proof\\off-chain, then pays\\the gas; reads the\\choice, can drop it};
  \node[svcbox, minimum width=2.8cm, minimum height=1.9cm] (registrar) at (7.0,-4.8)
    {\textbf{Registrar}\\checks eligibility,\\builds the tree\\off-chain; not in\\the ballot path};
  \node[svcbox, minimum width=2.8cm, minimum height=1.9cm] (indexer) at (10.4,-4.8)
    {\textbf{Indexer}\\tails events into a\\read cache,\\rebuildable\\from genesis};
  \node[svcbox, minimum width=2.8cm, minimum height=1.9cm] (pinning) at (13.8,-4.8)
    {\textbf{Metadata}\\pins the leaf set,\\manifestos and\\ballot descriptions\\to IPFS};
  \node[store, minimum width=2.8cm, minimum height=1.5cm] (pg) at (6.9,-7.6)
    {\textbf{PostgreSQL}\\voter records,\\commitments,\\indexed events};
  \node[store, minimum width=3.0cm, minimum height=1.5cm] (ipfs) at (12.0,-7.6)
    {\textbf{IPFS (Pinata)}\\published leaf set,\\ballot metadata};

  % ==== contract layer ====
  \node[cbox, minimum width=3.2cm, minimum height=1.6cm] (fwd) at (3.3,-12.05)
    {\texttt{VoteForwarder}\\ERC-2771: a\\convenience, not a\\security control};
  \node[cbox, minimum width=3.6cm, minimum height=1.6cm] (reg) at (3.7,-15.05)
    {\texttt{VoterRegistry}\\one root, roll size and\\leaf-set CID, frozen\\when registration closes;\\no commitment on-chain};
  \node[cbox, minimum width=4.0cm, minimum height=1.8cm] (elec) at (8.8,-13.55)
    {\texttt{Election}\\lifecycle state machine,\\nullifier set, ballot\\ceiling, tallies};
  \node[cbox, minimum width=3.8cm, minimum height=1.6cm] (ver) at (13.4,-12.05)
    {\texttt{SemaphoreVerifier}\\Groth16 membership\\proof; no Semaphore\\group contract on-chain};
  \node[cbox, minimum width=3.4cm, minimum height=1.4cm] (fac) at (13.4,-15.05)
    {\texttt{ElectionFactory}\\EIP-1167 minimal\\proxy clones};
  \node[cbox, minimum width=3.4cm, minimum height=1.2cm] (tally) at (8.8,-16.45)
    {\texttt{TallyLib}\\counting rules};

  % ==== layer bands ====
  \begin{scope}[on background layer]
    \node[layer, fit=(booth)(admin)(audit), inner sep=0.35cm] (bandA) {};
    \node[layer, fit=(relayer)(registrar)(indexer)(pinning)(pg)(ipfs),
          inner sep=0.35cm] (bandB) {};
    \node[layer, fit=(fwd)(fac)(elec)(ver)(reg)(tally), inner sep=0.35cm] (bandC) {};
  \end{scope}
  \node[layerlab, anchor=south east] at (bandA.north east)
    {Browser: Next.js 15, wagmi, viem};
  \node[layerlab, anchor=south west] at (bandB.north west)
    {Off-chain services: never a source of truth};
  \node[layerlab, anchor=south east] at (bandC.north east)
    {Contracts on Polygon Amoy (chain id 80002)};

  % ==== the sponsored ballot path ====
  \draw[ballot] ([xshift=0.6cm]booth.south) -- ([xshift=0.95cm]relayer.north);
  \node[keylab, anchor=north west] at (4.5,-1.85)
    {signed ballot and zk proof,\\as a meta-transaction};
  \draw[ballot] (relayer.south) -- ([xshift=0.1cm]fwd.north);
  \draw[ballot] (fwd.east) -- (elec.north west);

  % ==== the audit path: chain read, no backend ====
  \draw[audit] ([xshift=1.9cm]audit.south) -- ++(0,-12.25) -- (elec.east);
  \node[keylab, rotate=90, anchor=west] at (15.78,-9.6)
    {reads the chain directly, never our backend};
  \node[slab, anchor=south, text=blue!45!black] at (12.9,-13.43)
    {no backend in this path};

  % ==== escape hatch: voter submits their own transaction ====
  \draw[hatch] ([xshift=-1.6cm]booth.south) -- ++(0,-1.35) -- (0.7,-2.65)
    -- (0.7,-13.55) -- (elec.west);
  \node[lab, rotate=90, anchor=west, text=black!70] at (0.88,-10.0)
    {a voter holding gas can submit directly};

  % ==== ordinary traffic ====
  % the authority signs and funds its own transactions: the backend is not in
  % this path either
  \draw[flow] ([xshift=-0.5cm]admin.south) -- ([xshift=-0.1cm]elec.north);
  \node[lab, anchor=north west] at (8.95,-1.85)
    {the authority pays its own gas:\\create, open a phase, cancel};
  % the booth fetches the published leaf set itself and rebuilds the tree, so
  % the registrar never learns which leaf is voting
  \draw[flow, rounded corners=4pt] ([xshift=1.9cm]booth.south) -- ++(0,-1.55)
    -- (12.1,-2.85) -- ([xshift=0.1cm]ipfs.north);
  % the registrar's only transaction
  \draw[flow, rounded corners=4pt] (registrar.west) -- (5.2,-4.8)
    -- ([xshift=1.5cm]reg.north);
  \node[lab, anchor=north west, text width=3.0cm] at (5.4,-9.1)
    {one transaction: the root, the roll size and the leaf-set CID};
  \draw[flow] (registrar.south) -- (pg.north);
  \draw[flow] ([xshift=-0.6cm]indexer.south) -- (pg.north east);
  \draw[flow] (pinning.south) -- ([xshift=0.9cm]ipfs.north);
  \draw[flow] ([xshift=0.6cm]elec.north) -- ([xshift=-1.0cm]indexer.south);
  \node[lab, anchor=north west] at (9.6,-9.1)
    {\texttt{VoteCast} and\\\texttt{Finalized} logs};

  % ==== inside the contract layer ====
  \draw[flow] (fac.west) -- (elec.south east);
  \node[slab, anchor=north, rotate=28] at (11.3,-14.82) {clones};
  \draw[flow] (elec.north east) -- (ver.west);
  \node[slab, anchor=north, rotate=40] at (11.2,-12.32) {verifies};
  \draw[flow, rounded corners=4pt] ([xshift=-1.4cm]elec.south) -- (7.4,-15.05)
    -- (reg.east);
  \node[slab, anchor=south] at (6.6,-14.98) {reads the root};
  \draw[flow] (elec.south) -- (tally.north);
  \node[slab, anchor=west] at (8.95,-15.1) {counts};

\end{tikzpicture}

  \caption{\projectname{} component architecture.}
  \label{fig:architecture}
\end{figure}

Casting one ballot goes through every layer:

\begin{enumerate}
  \item The booth loads the voter's secret from browser storage, reads the
    frozen root and the list's IPFS address from the chain, downloads the list,
    and rebuilds the tree locally. The secret never leaves the device, and the
    registrar is not contacted.
  \item A background thread (Web Worker) builds the proof, the nullifier and
    the encoded candidate choice. Proving is the slowest step; B2 measures it
    on real phones.
  \item The booth signs the ballot request with a throwaway key (the signature
    moves no money) and sends it to the relayer. The relayer checks the proof
    before paying any gas, then submits it through \texttt{VoteForwarder}.
    Eligibility comes from the proof alone, so the signing key does not matter
    and the voter needs no wallet.
  \item \texttt{castVote} only runs during the \texttt{Voting} phase. It checks
    the proof against the frozen root and rejects a used nullifier or a ballot
    beyond the voter count. Otherwise it records the nullifier, adds one to the
    tally and emits a \texttt{VoteCast} event for the indexer.
\end{enumerate}

If any check fails, the whole transaction is reverted, so a ballot is either
counted or has no effect. The audit dashboard reads the chain through a public
RPC endpoint, never our backend: the registrar and relayer are conveniences
and the indexer is only a cache. Even if every one of our services is down or
lying, an observer can re-check every proof and recount the tally from the
\texttt{VoteCast} events. So nobody has to trust us for a correct result, but
voters still depend on us to register and, without gas, to vote.

Neither dependency is total: the registrar is needed only to register, and the
relayer only for voters with no gas. Because the relayer sees each choice and
IP address, it must run separately from the registrar, with its own database
and logs; this is a deployment requirement. Its ability to drop ballots is
limited, not removed (Section~\ref{sec:trust}).

\subsection{Smart contract design}

\begin{itemize}
  \item \texttt{ElectionFactory} is the only way to create an election. It
        deploys each one as a cheap EIP-1167 clone of a single
        \texttt{Election} contract, and keeps a list of them for auditors.
  \item \texttt{Election} holds the phase, the candidates, the used
        nullifiers and the tally. It treats a ballot sent through the
        forwarder exactly like a direct one. Proofs and nullifiers are tied
        to the election's ID, so they cannot be reused in another election.
        It calls a standalone Semaphore verifier with its own root.
  \item \texttt{VoterRegistry} holds the root, the voter count and the list's
        IPFS address, written in one registrar transaction. The root is frozen
        when registration closes. Capping ballots at the voter count costs one
        comparison.
  \item \texttt{TallyLib} holds the counting rules, so a new method can be
        added without touching the rest: plurality first, approval voting if
        time allows.
  \item \texttt{VoteForwarder} is OpenZeppelin's \texttt{ERC2771Forwarder}. It
        checks the signature, deadline and nonce, and gives relayed and direct
        ballots the same code path. It is a convenience, not a security
        measure: the proof stops a relayer changing a ballot, and the nullifier
        stops a replay. Each ballot uses a fresh signing key, because reusing
        one would link a voter's ballots across elections.
\end{itemize}

\begin{figure}[htbp]
  \centering
  % DVote election lifecycle state machine.
% Bare tikzpicture: \input this inside a figure environment.
% The preamble supplies tikz with arrows.meta, positioning, shapes.geometric,
% fit, backgrounds and calc.
%
% Greyscale-safe conventions:
%   single outline = a phase the election passes through
%   double outline = a terminal phase reached by counting a result
%   dashed outline = the terminal phase reached by abandoning the election
%
% Each transition carries its call, its time gate and who may make it. The
% authority holds only configuration and cancellation, because a transition
% only it can make is also a transition it can withhold.
\begin{tikzpicture}[
    font=\footnotesize,
    phase/.style    = {draw=black, thick, rounded corners=3pt, fill=white,
                       align=center, inner sep=3pt,
                       minimum width=4.6cm, minimum height=1.0cm},
    terminal/.style = {phase, double, double distance=1.1pt},
    aborted/.style  = {phase, dashed, fill=black!7,
                       minimum width=3.0cm},
    trans/.style     = {-{Latex[length=2.6mm,width=2.2mm]}, draw=black,
                       line width=1.1pt},
    abort/.style    = {-{Latex[length=2.2mm,width=1.8mm]}, draw=black!65,
                       semithick, dash pattern=on 2.5pt off 2pt,
                       rounded corners=3pt},
    tlab/.style     = {align=left, inner sep=2pt, text width=6.9cm,
                       font=\footnotesize,
                       execute at begin node={\hyphenpenalty=10000\exhyphenpenalty=10000}},
  ]

  % ==== the phases ====
  \node[phase]    (created)  at (6.0,0)
    {\texttt{Created}\\cloned by the factory,\\candidates and windows set};
  \node[phase]    (regn)     at (6.0,-2.8)
    {\texttt{Registration}\\identity commitments collected};
  \node[phase]    (voting)   at (6.0,-5.6)
    {\texttt{Voting}\\root frozen, ballots accepted};
  \node[phase]    (tallying) at (6.0,-8.4)
    {\texttt{Tallying}\\ballots refused; the counts\\already stand in storage};
  \node[terminal] (final)    at (6.0,-11.2)
    {\texttt{Finalized}\\result written, immutable};
  \node[aborted]  (cancelled) at (1.5,-1.4)
    {\texttt{Cancelled}\\terminal, no result};

  % ==== the forward path ====
  \draw[trans] (created.south)  -- (regn.north);
  \draw[trans] (regn.south)     -- (voting.north);
  \draw[trans] (voting.south)   -- (tallying.north);
  \draw[trans] (tallying.south) -- (final.north);

  \node[tlab, anchor=west] at (8.6,-1.4)
    {\texttt{openRegistration()}, time-gated: \texttt{block.timestamp} must have
     reached the announced opening, after which anyone may call it. A transition
     only the authority can make is one the authority can also withhold.};
  \node[tlab, anchor=west] at (8.6,-4.2)
    {\texttt{openVoting()}, time-gated on the voting window and likewise callable
     by anyone. It also requires a non-empty leaf-set CID beside the root, so a
     registrar that withholds the published roll stops the election rather than
     going unnoticed. The root was written when registration closed and cannot
     change from here.};
  \node[tlab, anchor=west] at (8.6,-7.0)
    {\texttt{closeVoting()}, time-gated on the voting window closing, and
     callable by anyone once that timestamp has passed, so the authority
     cannot stop the count.};
  \node[tlab, anchor=west] at (8.6,-9.8)
    {\texttt{finalize()}, likewise callable by anyone, the authority
     included. Writes the result and emits \texttt{Finalized}; no later call
     can change it.};

  % ==== the one abandonment path ====
  \draw[abort] (created.west) -- (3.4,0) -- (3.4,-0.9) -- (cancelled.north east);
  \draw[abort] (regn.west)    -- (3.4,-2.8) -- (3.4,-1.9) -- (cancelled.south east);
  \node[align=left, anchor=north, text width=3.3cm, inner sep=2pt,
        execute at begin node={\hyphenpenalty=10000\exhyphenpenalty=10000}] at (1.5,-2.2)
    {\texttt{cancel()}, the authority only, and only before voting opens. After
     that a ballot may exist, and cancelling the election would be a veto over
     a result the authority can already see.};

\end{tikzpicture}

  \caption{Election lifecycle. Every transition needs an explicit call after
    its scheduled time, and no transition goes back to an earlier phase.}
  \label{fig:lifecycle}
\end{figure}

An election starts as \texttt{Created}, is configured by the Election
Authority, and moves through \texttt{Registration}, \texttt{Voting} and
\texttt{Tallying} to \texttt{Finalized} (Figure~\ref{fig:lifecycle}). Only the
authority can cancel it, and only before voting opens.

Each phase change needs two things: the scheduled time must have passed, and
someone must send a transaction. The time check means no phase starts or ends
early. The transaction means every change is recorded on-chain with a sender.
Time alone would leave the current phase unclear; transactions alone would let
the authority close voting as soon as the count turned against it. Each step
only moves forward, so a phase cannot be rewound.

Only the authority can configure or cancel an election, since those are
administrative decisions. Every other step (opening registration, opening and
closing voting, publishing the result) can be done by anyone once its time
comes. An authority that dislikes the result gains nothing by refusing to act,
and cannot change the roll because the root is already frozen.

\subsection{Anonymous eligibility and double-vote prevention}

The system uses Semaphore v4. The voter's secret stays on their device, and the
registrar records only its commitment as a leaf of the Merkle tree. The browser
proves that it knows the secret behind one of the leaves, without revealing
which one. The nullifier is a hash of the secret and the election ID. It is
always the same for the same voter and election, so a second ballot is
rejected; it cannot be reversed to find the secret; and it differs between
elections, so a voter's ballots in two elections cannot be linked. The
candidate choice is locked into the proof, so a relayer cannot change it.

A voter hides among everyone on the roll, no more. The choice itself is public:
it is in the ballot, and the tally updates in the same transaction. What is
hidden is who cast it. Ballots are timestamped on-chain, so in a small election
it is easier to guess, and the very first ballot is not hidden at all. A demo
must therefore cast several ballots before any live one.

Groth16 needs a trusted setup. We use Semaphore's published setup files and
check their hashes against the Semaphore repository; the remaining assumption is
in Section~\ref{sec:trust}.

The simple alternative is to record \texttt{hasVoted[msg.sender]}. It is cheaper
but unacceptable: each ballot would show the voter's address next to their
choice, and the registrar knows which person owns which address, so it would
publish how everyone voted and make coercion easy to enforce. The nullifier
costs the same storage write as the flag, plus the proof check, so anonymity
has a price; A1 measures it.

\subsection{Scalability strategy}

We describe the system in three tiers (Table~\ref{tab:tiers}), and build only
the last two. Tier 1 is only for comparison: its registration cost grows with
the number of voters, and its ballots show who voted.

\begin{table}[htbp]
  \centering
  \caption{Build tiers.}
  \label{tab:tiers}
  \small
  \begin{tabular}{@{}lL{3.9cm}L{5.3cm}L{3.0cm}@{}}
    \toprule
    \textbf{Tier} & \textbf{Eligibility} & \textbf{Ballot privacy}           & \textbf{Voter pays gas} \\
    \midrule
    1             & On-chain roll        & None; ballots linked to addresses & Yes                     \\
    2             & Merkle root          & None                              & Yes                     \\
    3             & Zero-knowledge proof & Anonymous, nullifier-protected    & No                      \\
    \bottomrule
  \end{tabular}
\end{table}

Two operations cost the same gas at any roll size: publishing the roll (one
transaction writing the root, the voter count and the IPFS address), and
casting a tier 3 ballot (one proof check of fixed size). Both come from building
the tree off-chain and putting only its root on-chain. We do not claim that
nothing grows:

\begin{itemize}
  \item Total election gas grows with the number of ballots cast, since each
        ballot is a transaction.
  \item Browser work grows: rebuilding the tree grows with the roll, and
        proving grows slightly with the tree's depth. C2 measures both.
  \item Off-chain work grows with the roll: the registrar builds the tree, and
        each auditor rebuilds it. Hashing off-chain is still far cheaper than
        storing on-chain.
\end{itemize}

We do not use Semaphore's own group contracts for this reason. They store the
group on-chain and update the root for every new member, which is simpler but
costs one transaction per voter. We use Semaphore's circuits, verifier and
nullifiers, but not its on-chain groups. Tier 2 already publishes only a root,
so zero knowledge changes the ballot, not registration.

In tier 2 the ballot sends its Merkle proof to the contract, costing a few
hundred gas per tree level. In tier 3 that proof stays private inside the
zero-knowledge proof, so the ballot cost stops depending on the tree's depth.
We expect tier 3 to be flat but always more expensive than tier 2, because the
proof check costs far more than the path it replaces. Anonymity costs gas; it
does not save it.

With only a root on-chain, nobody can list the voters, so the registrar
publishes the list on IPFS and writes its address next to the root. An auditor
rebuilds the root from the list and catches any mismatch. That does not show
whose commitments are on the list, so before voting the booth checks that the
voter's own commitment is there; a voter who was swapped out will notice. There
is a gap between the root freezing and voting opening so that voters can act on
this, and voting cannot open until the list's address is published. None of
this makes the roll honest (Section~\ref{sec:trust}).

C1 measures registration and ballot gas for rolls of $10$ to $10^5$ voters.
Flat registration follows from the design, but C1 checks the actual code: a
contract that stored commitments would show a rising line. If a line is not
flat, Evaluation~2 will say so. C2 covers the other risk: whether a
hundred-thousand-voter roll can be handled in a browser at all.

\subsection{Technology choices and justification}

Table~\ref{tab:stack} lists each choice, the main alternative, and the reason.

\begin{table}[htbp]
  \centering
  \caption{Technology choices.}
  \label{tab:stack}
  \small
  \begin{tabular}{@{}L{2.5cm}L{3.2cm}L{3.4cm}L{7.2cm}@{}}
    \toprule
    \textbf{Layer}      & \textbf{Choice}             & \textbf{Alternative}           & \textbf{Reason}                                           \\
    \midrule
    Contract language   & Solidity 0.8.28             & Vyper, non-EVM chains          & Largest ecosystem; used in the course; one language throughout\\
    Contract toolchain  & Foundry                     & Hardhat 3                      & Tests in Solidity; \texttt{forge snapshot} gives A1 directly \\
    Contract libraries  & OpenZeppelin 5.x            & Writing our own                & Audited access control, forwarder, clones and Merkle proofs            \\
    Election deployment & EIP-1167 clones             & Full deploy per election       & Much cheaper; cannot be upgraded during a vote            \\
    Anonymity           & Semaphore v4                & Custom circuits, MACI          & Audited proofs and nullifiers; custom circuits too costly to build        \\
    Gasless voting      & ERC-2771 forwarder          & Voter pays gas                 & Voter needs no cryptocurrency; direct submission as fallback    \\
    Testbed             & Ethereum Sepolia            & Mainnet, other testnets        & Free test ETH, behaves like mainnet; the recommended testnet for contract development \\
    Frontend            & Next.js, wagmi, viem        & Plain SPA, ethers v6           & Typed contract bindings; can be exported as a static site               \\
    Backend             & Fastify, PostgreSQL, BullMQ & Express, MongoDB               & Relational data for the indexer, shared types, reliable job queue        \\
    Metadata storage    & IPFS via Pinata             & Storing on-chain               & Same tamper-evidence for far less gas                           \\
    \bottomrule
  \end{tabular}
\end{table}

A few choices need more explanation:

\begin{itemize}
  \item \emph{Solidity and Foundry.} Vyper has few libraries, and a non-EVM
        chain would lose the ability to run the same code locally and on the
        testnet, which is what lets us measure gas locally. Foundry's
        \texttt{forge snapshot} gives A1 and C1, and its fuzzer gives the
        attack tests D1 needs. Hardhat 3 integrates better with TypeScript,
        but wagmi generates types from Foundry's output anyway.
  \item \emph{Two different trees.} OpenZeppelin's \texttt{MerkleProof} is only
        used in tier 2. Tier 3 uses a different hash, so the tree builder, the
        root and the verification are replaced. The phases, factory, indexer
        and gas baseline carry over.
  \item \emph{Clones.} An EIP-1167 clone is a tiny contract that forwards to one
        shared implementation, so each election is cheap to deploy. A bug in
        that implementation would affect every election, but it cannot be
        changed, whereas an upgradeable contract would let the authority change
        the rules mid-vote.
  \item \emph{Semaphore v4.} It provides exactly what we need, anonymous
        membership with a per-election nullifier, and it is audited. Writing our
        own circuits would take too long and put unaudited cryptography at the
        centre of a security project. MACI would add coercion resistance, but
        needs a trusted coordinator, encrypted messages and key changes.
  \item \emph{IPFS.} An IPFS address guarantees the content has not changed,
        but not that it is still available: content that nobody pins
        disappears. That is acceptable for a manifesto, and it is why both we
        and the auditors pin the voter list.
\end{itemize}

\subsection{Deployment testbed}

Everything is deployed to Ethereum Sepolia, the Ethereum Foundation's
recommended testnet for application and contract development. It has the best
faucet and tooling support of any testnet, which matters when we need to fund
many test ballots, and it runs the same way as Ethereum mainnet.

Gas is deterministic: the same code and input use the same gas on any chain. So
the gas parameters (A1, C1) are measured on a local Foundry node, where testing
roll sizes from ten to a hundred thousand is instant and free. The load and
reorg tests also run locally, so that network noise does not distort them.
Only the end-to-end latency test and the live demo need Sepolia.
